% =====================================================================
% Canvas3D Figure 5: ArtiFlow - Variants A/B/C  (LAYOUT v3, collision-free)
%
% Geometry summary (all three variants share the SAME layout):
%
%       [sc]x --> [rend] --> [fr] --> [vd]  --> [out] (down)
%                  |    \--> [df]
%                  |    \--> [ef]
%                  v
%                 [cn]  (a full row below rend, x-aligned with rend)
%
% Key invariants kept by construction:
%   * vd sits on the TOP row (right of fr), never in the same vertical
%     band as df/ef/cn.
%   * cn sits on the BOTTOM row directly below rend, with a clear
%     vertical gutter between the rend column and the df/ef column.
%   * df and ef bleed down and left into cn through that gutter,
%     without touching rend, vd, or each other.
%   * fr attaches to vd via a straight east-west arrow; no big arcs.
% =====================================================================

% ============================ VARIANT A ==============================
\begin{figure}[t]
    \centering
    \resizebox{\linewidth}{!}{%
    \begin{tikzpicture}[
        font=\sffamily,
        bx/.style={draw=none, rounded corners=3pt, minimum width=2.4cm,
                   minimum height=0.72cm, align=center,
                   font=\sffamily\scriptsize, thick, text=white},
        data/.style={draw=black!25, rounded corners=2pt, fill=white,
                     font=\sffamily\tiny, minimum width=2.1cm,
                     minimum height=0.42cm, align=center},
        ar/.style={->, >=Stealth, very thick, draw=black!45}
    ]
        % --- Input ---
        \node[data, fill=yellow!15, draw=yellow!50!black] (sc)
            {$\mathcal{S}^*$, $\mathcal{D}^*$};

        % --- Renderer ---
        \node[bx, fill=red!65, right=1.1cm of sc,
              minimum width=2.5cm, minimum height=1.05cm] (rend)
            {\textbf{Multi-View Render}\\{\tiny 3DGS rasterize}\\{\tiny $I_t, D_t, E_t$}};

        % --- Render products stack on rend's right ---
        \node[data, right=1.2cm of rend, minimum width=1.8cm] (fr) {RGB $I_t$};
        \node[data, below=0.45cm of fr,   minimum width=1.8cm] (df) {Depth $D_t$};
        \node[data, below=0.45cm of df,   minimum width=1.8cm] (ef) {Edge $E_t$};

        % --- ControlNet sits a FULL row below rend, x-aligned. ---
        \node[bx, fill=red!45, below=2.2cm of rend,
              minimum width=2.5cm, minimum height=1.05cm] (cn)
            {\textbf{Dual ControlNet}\\{\tiny Depth $s_d=0.85$}\\{\tiny HED $s_e=0.6$}};

        % --- Video DiT on the top row (right of fr), same baseline as rend. ---
        \node[bx, fill=purple!60, right=1.2cm of fr,
              minimum width=2.6cm, minimum height=1.05cm] (vd)
            {\textbf{Video DiT}\\{\tiny Wan2.2 / LTX2.3}\\{\tiny temp.\ attention}};

        % --- Final output below vd, but first we place an invisible anchor
        %     half a row down so `out` sits clearly under vd. ---
        \node[data, below=2.2cm of vd, minimum width=2.3cm] (out)
            {Output Video $V$};

        % ================= EDGES =================
        % Top row: straight east-west arrows.
        \draw[ar] (sc.east)   -- (rend.west);
        \draw[ar] (rend.east) -- (fr.west);
        \draw[ar] (rend.east) to[out=0,in=180,looseness=1.0] (df.west);
        \draw[ar] (rend.east) to[out=0,in=180,looseness=1.2] (ef.west);
        \draw[ar] (fr.east)   -- (vd.west);

        % Straight vertical drop from rend to cn: nothing sits in between.
        \draw[ar] (rend.south) -- (cn.north);

        % df -> cn: down-left diagonal inside the empty gutter between the
        % rend column and the df column.
        \draw[ar] (df.west) to[out=180,in=60,looseness=1.0] (cn.north east);

        % ef -> cn: drop down in the gutter to the level of cn and come in
        % from the EAST, so it never crosses vd (vd is on the top row).
        \draw[ar] (ef.west) to[out=180,in=15,looseness=1.0] (cn.east);

        % cn -> vd: gentle rising diagonal into vd.south, through the empty
        % space between the cn row and the df/ef stack.
        \draw[ar] (cn.east) to[out=0,in=-90,looseness=1.0] (vd.south);

        % vd -> out : straight vertical drop (out is directly below vd).
        \draw[ar] (vd.south) -- (out.north);
    \end{tikzpicture}}
    \caption{\textbf{ArtiFlow (Variant A).} Flat style with red for rendering,
    darker red for ControlNet, purple for the diffusion model.}
    \label{fig:artiflow_A}
\end{figure}

% ============================ VARIANT B ==============================
\begin{figure}[t]
    \centering
    \resizebox{\linewidth}{!}{%
    \begin{tikzpicture}[
        font=\sffamily,
        bx/.style={draw=black!80, rounded corners=1pt, minimum width=2.5cm,
                   minimum height=0.72cm, align=center,
                   font=\sffamily\scriptsize, thick},
        data/.style={draw=black!40, rounded corners=1pt, fill=white,
                     font=\sffamily\tiny, minimum width=2.1cm,
                     minimum height=0.42cm, align=center},
        ar/.style={->, >=Stealth, thick, draw=black!70}
    ]
        \node[data] (sc) {$\mathcal{S}^*$, $\mathcal{D}^*$};

        \node[bx, right=1.1cm of sc, minimum width=2.6cm,
              minimum height=1.05cm] (rend)
            {Multi-View Render\\{\tiny 3DGS rasterize}\\{\tiny $I_t, D_t, E_t$}};

        \node[data, right=1.2cm of rend, minimum width=1.9cm] (fr) {RGB $I_t$};
        \node[data, below=0.45cm of fr,   minimum width=1.9cm] (df) {Depth $D_t$};
        \node[data, below=0.45cm of df,   minimum width=1.9cm] (ef) {Edge $E_t$};

        \node[bx, below=2.2cm of rend, minimum width=2.6cm,
              minimum height=1.05cm] (cn)
            {Dual ControlNet\\{\tiny Depth + HED}\\{\tiny $s_d=0.85$, $s_e=0.6$}};

        \node[bx, right=1.2cm of fr, minimum width=2.7cm,
              minimum height=1.05cm] (vd)
            {Video DiT\\{\tiny Wan2.2 / LTX2.3}\\{\tiny temporal attention}};

        \node[data, below=2.2cm of vd, minimum width=2.4cm] (out)
            {Output Video $V$};

        \draw[ar] (sc.east)   -- (rend.west);
        \draw[ar] (rend.east) -- (fr.west);
        \draw[ar] (rend.east) to[out=0,in=180,looseness=1.0] (df.west);
        \draw[ar] (rend.east) to[out=0,in=180,looseness=1.2] (ef.west);
        \draw[ar] (fr.east)   -- (vd.west);

        \draw[ar] (rend.south) -- (cn.north);
        \draw[ar] (df.west) to[out=180,in=60,looseness=1.0] (cn.north east);
        \draw[ar] (ef.west) to[out=180,in=15,looseness=1.0] (cn.east);

        \draw[ar] (cn.east) to[out=0,in=-90,looseness=1.0] (vd.south);
        \draw[ar] (vd.south) -- (out.north);
    \end{tikzpicture}}
    \caption{\textbf{ArtiFlow (Variant B).} Monochrome line style, no fills,
    maximum clarity.}
    \label{fig:artiflow_B}
\end{figure}

% ============================ VARIANT C ==============================
\begin{figure}[t]
    \centering
    \resizebox{\linewidth}{!}{%
    \begin{tikzpicture}[
        font=\sffamily,
        bx/.style={draw=black!60, fill=black!8, rounded corners=2pt,
                   minimum width=2.5cm, minimum height=0.72cm, align=center,
                   font=\sffamily\scriptsize, thick},
        acc/.style={draw=blue!60!black, fill=blue!12, rounded corners=2pt,
                    minimum width=2.5cm, minimum height=0.72cm, align=center,
                    font=\sffamily\scriptsize, thick},
        data/.style={draw=black!35, fill=white, rounded corners=1pt,
                     font=\sffamily\tiny, minimum width=2.1cm,
                     minimum height=0.42cm, align=center},
        ar/.style={->, >=Stealth, thick, draw=black!55}
    ]
        \node[data] (sc) {$\mathcal{S}^*$, $\mathcal{D}^*$};

        \node[bx, right=1.1cm of sc, minimum width=2.6cm,
              minimum height=1.05cm] (rend)
            {Multi-View Render\\{\tiny 3DGS rasterize}\\{\tiny $I_t, D_t, E_t$}};

        \node[data, right=1.2cm of rend, minimum width=1.9cm] (fr) {RGB $I_t$};
        \node[data, below=0.45cm of fr,   minimum width=1.9cm] (df) {Depth $D_t$};
        \node[data, below=0.45cm of df,   minimum width=1.9cm] (ef) {Edge $E_t$};

        \node[acc, below=2.2cm of rend, minimum width=2.6cm,
              minimum height=1.05cm] (cn)
            {\textbf{Dual ControlNet}\\{\tiny Depth + HED}\\{\tiny $s_d=0.85$, $s_e=0.6$}};

        \node[bx, right=1.2cm of fr, minimum width=2.7cm,
              minimum height=1.05cm] (vd)
            {Video DiT\\{\tiny Wan2.2 / LTX2.3}\\{\tiny temporal attention}};

        \node[data, below=2.2cm of vd, minimum width=2.4cm] (out)
            {Output Video $V$};

        \draw[ar] (sc.east)   -- (rend.west);
        \draw[ar] (rend.east) -- (fr.west);
        \draw[ar] (rend.east) to[out=0,in=180,looseness=1.0] (df.west);
        \draw[ar] (rend.east) to[out=0,in=180,looseness=1.2] (ef.west);
        \draw[ar] (fr.east)   -- (vd.west);

        \draw[ar] (rend.south) -- (cn.north);
        \draw[ar] (df.west) to[out=180,in=60,looseness=1.0] (cn.north east);
        \draw[ar] (ef.west) to[out=180,in=15,looseness=1.0] (cn.east);

        \draw[ar] (cn.east) to[out=0,in=-90,looseness=1.0] (vd.south);
        \draw[ar] (vd.south) -- (out.north);
    \end{tikzpicture}}
    \caption{\textbf{ArtiFlow (Variant C).} Grayscale with blue accent on the
    dual-ControlNet conditioning, highlighting the control mechanism.}
    \label{fig:artiflow_C}
\end{figure}
